%!TEX root = ../index.tex
% - - - - - - - - - - - - - - - - - - - - - - - - - - - - - - - - - - - -

\thispagestyle{empty}
{
	\sffamily\scriptsize
	\setlength\parindent{0em}
	
	\begin{wrapfigure}[2]{r}[0.0cm]{2cm}
		\vspace{-0.75cm}
		\includegraphics[width=2.0cm]{fcw}
	\end{wrapfigure}
	
	\textbf{(CC):BY-SA 4.0}
	
	\vspace{.5\baselineskip}
	
	Textos publicados bajo licencia CREATIVE COMMONS 4.0 internacional, que permite:
	
	\begin{itemize}
	\item \textbf{Compartir} // copiar y redistribuir el material en cualquier medio o formato.
	\item \textbf{Adaptar} // remezclar, transformar y crear a partir del material.
	\end{itemize}
	
	Bajo las condiciones siguientes:
	
	\begin{itemize}
	\item \textbf{Reconocimiento} // Debes reconocer adecuadamente la autoría, proporcionar un enlace a la licencia e indicar si se han realizado cambios. Puedes hacerlo de cualquier manera razonable, pero no de una manera que sugiera que tienes nuestro apoyo o lo recibes por el uso que haces.
	\item \textbf{Compartir igual} // Si remezclas, transformas o creas a partir del material, deberás difundir sus contribuciones bajo la misma licencia que el original.
	\end{itemize}
	
	No hay restricciones adicionales. No puedes aplicar términos legales o medidas tecnológicas que legalmente restrinjan realizar aquello que la licencia permite.
	
	\vspace{1.0\baselineskip}
	
	\textbf{Aviso}: No tienes que cumplir con la licencia para aquellos elementos del material en el dominio público o cuando su utilización esté permitida por la aplicación de una excepción o un límite. No se dan garantías. La licencia puede no ofrecer todos los permisos necesarios para la utilización prevista. Por ejemplo, otros derechos como los de publicidad, privacidad, o los derechos morales pueden limitar el uso del material.
	
	\vspace{1.0\baselineskip}
	
	Para más información puedes consultar:\par
	
	\vspace{0.5\baselineskip}
	
	{\setlength\parindent{1cm}
	> ~ \url{https://creativecommons.org/licenses/by-sa/4.0/}
	}
	
	\vspace{1.0\baselineskip}
	
	Este fanzine se distribuye en papel en:
	\begin{linelist}
	\item Un lugar;
	\item Otro lugar;
	\item Otro lugar;
	\item Otro lugar;
	\item Otro lugar;
	\item Otro lugar
	\end{linelist}.
	
	\vspace{1.0\baselineskip}
	
	A Tamagotchi lo puedes encontrar en:\par
	
	\vspace{0.5\baselineskip}
	
	{\setlength\parindent{1cm}
	> ~ \url{http://wearetamagotchi.com}
	}

}
\rmfamily
\clearpage